%! TeX program = lualatex
%! TEX options = -synctex=1 -interaction=nonstopmode -file-line-error --shell-escape "%DOC%"
\documentclass[10pt,a5paper]{mszalik}

\begin{document}

\thispagestyle{empty}

\begin{center}
	\vspace*{0.5cm}


	\bfseries\scshape
	\huge NIESZPORY LUDOWE

    \vspace*{0.5cm}

    \huge NIEDZIELA CHRYSTUSA KRÓLA

    \vspace*{0.5cm}

	\vfill

	\begin{figure}[h]
		\centering
		\includegraphics[width=0.45\linewidth]{logo.png}
	\end{figure}

	\vfill

    \vspace*{0.5cm}

    % {\large \textbf{25 października 2026}}

    % \bigskip

    {\large Kolegiata Świętego Krzyża, Wrocław}

\end{center}

\newpage

\vspace*{2cm}

\vspace*{\fill}

\hrulefill

{\footnotesize

    Śpiewy i teksty liturgiczne oraz ich tłumaczenia pochodzą z:

	\begin{itemize}[leftmargin=*]
		\item \textit{ks. Bronisław Gładysz -- Hymny brewjarza rzymskiego oraz patronału polskiego: przekład i objaśnienia} (1933),
        \item \textit{o. Gaspar Lefebvre -- Mszał Rzymski z dodaniem Nabożeństw Nieszpornych} (1949),
        \item \textit{Liber Usualis} (1961),
        \item \textit{ks. Antoni Reginek -- Wybór pieśni nabożnych Franciszka Karpińskiego oraz psalmów w jego tłumaczeniu} (2001),
        \item \textit{ks. Jan Siedlecki -- Śpiewnik kościelny, wydanie XLI} (2021),
	\end{itemize}

	W niektórych przypadkach wykorzystano cyfrową wersję tekstów
	opublikowaną na portalach \textit{sbc.org.pl}, \textit{divinumofficium.com} i \textit{missalemeum.com}.

    \medskip

    Zapis gabc neumów: Piotr Matuszak; Andrew Hinkley, Felipe Gasper \textit{gregobase.selapa.net}.

	\medskip

    Duszpasterstwo Wiernych Tradycji Łacińskiej w Archidiecezji Wrocławskiej jest duszpasterstwem
    erygowanym przez J.E. ks. abp. Józefa Kupnego, Metropolitę Wrocławskiego.

	\bigskip

	Opracowanie: Mateusz Burszewski, Piotr Matuszak\\
	Skład i łamanie: Michał Siemaszko, Piotr Matuszak\\

	\texttt{tradycja.archidiecezja.wroc.pl}
}


%%%%%%%%%%%%%%%%%%%%%%%%%%%%%%%%%%%%%%%%%%%%%%%%%%%%%%%%%%%%%%%%%%%%%%%%%%

\newpage

\fancyhead[CE]{II. VESPERÆ}
\fancyhead[CO]{DOMINI NOSTRI IESU CHRISTI REGIS}

\section{NIESZPORY LUDOWE}

\rubric{Kapłan i najbliższa asysta przyklękają przed ołtarzem. Klękają i modlą się w ciszy po czym wstają i przechodzą do sedili.}


\gresetinitiallines{0}
\gregorioscore{Melody/boze-wejrzyj.gabc}
\gresetinitiallines{1}

\medskip

Chwała bądź Bogu w Trójcy \textbf{je}dynemu: \redcol{*} Ojcu, Synowi, Ducho\textit{wi} Świętemu. \\
Jak od początku była, \textbf{tak} i ninie, \redcol{*} I na wiek wieków niechaj \textit{zaw}sze słynie. \\
Alleluja.

\proroctwo{Antyfona}

\gregorioscore{Melody/an--pacificus_vocabitur--solesmes_1961.1.gabc}

\medskip
Książeciem pokoju nazwany będzie, a stolica Jego będzie umocniona na wieki.

\proroctwo{Psalm 109}

\begin{lilypond}
\relative c' {
    \key f \major
    \clef treble
    \override Staff.TimeSignature.stencil = ##f
    \override Staff.NoteHead.style = #'altdefault
    \cadenzaOn
    a'\breve bes8 a8 \parenthesize a8 g8 a4 \bar "|" \break
    g\breve a8 \parenthesize a8 g8 f8 bes8([a8]) g4 \bar "||"
    \cadenzaOff
}
\addlyrics {
    \override LyricText.font-size = #-1
    Rzekł_Pan_do_Pana_ \markup { \bold me } -- go łas -- ka -- wym,
    Swym_głosem:_siądź_mi \markup { \bold przy } _ bo -- ku pra -- wym.
}
\end{lilypond}

% 1. Rzekł Pan do Pana \textbf{me}go łaskawym, \redcol{*} Swym głosem: siądź mi \textbf{przy} boku prawym.\\
2. Aż Twoje wszystkie \textbf{zu}chwałe wrogi, \redcol{*} Dam za podnóżek \textbf{pod} Twoje nogi.\\
3. Berło Twej mocy \textbf{wy}dam z Syjonu, \redcol{*} świat cały padnie \textbf{u} Twego tronu.\\
4. A Ty używać \textbf{bę}dziesz praw swoich, \redcol{*} wśród nieprzyjaciół \textbf{pa}nując Twoich.\\
5. Skoro rozpoczniesz \textbf{swe} panowanie, \redcol{*} przy Tobie lud Twój w \textbf{jas}ności stanie.\\
6. Nim jeszcze zorzy \textbf{świe}cić kazałem, \redcol{*} Ciebie z wnętrzności \textbf{swo}ich wydałem.\\
7. Pan to poprzysiągł, \textbf{Jego} zaś mowa \redcol{*} danego nigdy \textbf{nie} cofnie słowa.\\
8. Ty jesteś kapłan \textbf{do} końca wieka, \redcol{*} według obrządku \textbf{Mel}chizedeka.\\
9. Przy Twej prawicy \textbf{Pan} jest nad pany, \redcol{*} w dzień gniewu swego \textbf{ze}trze tyrany.\\
10. Sąd swój rozciągnie \textbf{po} całym świecie \redcol{*} i nieposłuszne \textbf{na}rody zgniecie.\\
11. Pyszną na ziemi \textbf{gło}wę poniży, \redcol{*} która Mu Jego \textbf{chwa}ły ubliży.\\
12. Z mętnej po drodze \textbf{pić} będzie rzeki, \redcol{*} dlatego głowę \textbf{wznie}sie na wieki.\\
13. Chwała bądź Bogu w Trójcy \textbf{Je}dynemu, \redcol{*} Ojcu, Synowi, Du\textbf{cho}wi Świętemu.\\
14. Jak od początku była, \textbf{tak} i ninie, \redcol{*} I na wiek wieków \textbf{nie}chaj zawsze słynie.\\

\proroctwo{Antyfona}

\gregorioscore{Melody/an--pacificus_vocabitur--solesmes_1961.1.gabc}

\medskip
Książeciem pokoju nazwany będzie, a stolica Jego będzie umocniona na wieki.

\proroctwo{Antyfona}

\gregorioscore{Melody/an--regnum_ejus--solesmes_1961.1.gabc}

\medskip
Królestwo Jego Królestwo wieczne jest i wszystkie narodu służyć Mu będą i słuchać Go będą.

\newpage
\proroctwo{Psalm 110}

\begin{lilypond}
\relative c' {
    \key f \major
    \clef treble
    \override Staff.TimeSignature.stencil = ##f
    \override Staff.NoteHead.style = #'altdefault
    \cadenzaOn
    f\breve d8 d4 c4 \bar "|" \break
    f\breve g8 a8 g4 f4 \bar "||"
    \cadenzaOff
}
\addlyrics {
    \override LyricText.font-size = #-1
    Całym_Cię_sercem_chwalić_bę -- \markup { \bold dą, } Pa -- nie,
    gdzie_rada_mędrców,_gdzie_ca -- \markup { \bold ły } zbór sta -- nie.
}
\end{lilypond}

% 1. Całym Cię sercem chwalić bę\textbf{dę}, Panie, gdzie rada mędrców, gdzie ca\textbf{ły} zbór stanie.
2. Wielkość niezmierna dzieł Boga \textbf{Ja}kuba, \redcol{*} Którymi władnie, jak Mu \textbf{się} podoba.\\
3. Co pocznie, wszystko pełne jest \textbf{zac}ności, \redcol{*} Pełne uwielbień. Wiek spra\textbf{wie}dliwości.\\
4. Jego nie przetrwa żaden, a \textbf{na} ziemi, \redcol{*} Wieczną pamiątkę sprawi \textbf{cu}dy Swymi.\\
5. Pan dobrotliwy, Pan to mi\textbf{ło}sierny, \redcol{*} Karmi i hojnie boga\textbf{ci} lud wierny.\\
6. Stateczny w słowie, co raz po\textbf{sta}nowił, \redcol{*} Wiecznymi czasy tego \textbf{nie} odmówił.\\
7. Moc Swą okazał, gdy wygnał \textbf{po}gany, \redcol{*} A ten kraj synom Izra\textbf{e}la dany.\\
8. Sprawy rąk Jego prawdą miar\textbf{ko}wane, \redcol{*} A zaś słusznością prawa \textbf{pro}stowane.\\
9. Te się na wieki żadnych lat \textbf{nie} boją, \redcol{*} Bo na słuszności i na \textbf{praw}dzie stoją.\\
10. Pan sługi Swoje z niewoli \textbf{wy}bawił, \redcol{*} I dla nich wieczny testa\textbf{ment} zostawił.\\
11. Imię ma straszne i pełne \textbf{świę}tości, \redcol{*} A bojaźń Boża począt\textbf{kiem} mądrości.\\
12. Ci, co z niej biorą wzór życia \textbf{na} ziemi, \redcol{*} Będą od wszystkich wiecznie \textbf{chwa}lonymi.\\
13. Chwała bądź Bogu w Trójcy Je\textbf{dy}nemu, \redcol{*} Ojcu, Synowi, Ducho\textbf{wi} Świętemu.\\
14. Jak od początku była, tak \textbf{i} ninie, \redcol{*} I na wiek wieków niechaj \textbf{zaw}sze słynie.\\

\proroctwo{Antyfona}

\gregorioscore{Melody/an--regnum_ejus--solesmes_1961.1.gabc}

\medskip
Królestwo Jego Królestwo wieczne jest i wszystkie narodu służyć Mu będą i słuchać Go będą.

\newpage
\proroctwo{Antyfona}

\gregorioscore{Melody/an--ecce_vir_oriens--solesmes_1961.1.gabc}

\medskip
Oto mąż, a Wschód imię Jego, zasiądzie i zapanuje i pokój narodom opowie.

\proroctwo{Psalm 111}

\begin{lilypond}
\relative c' {
    \key f \major
    \clef treble
    \override Staff.TimeSignature.stencil = ##f
    \override Staff.NoteHead.style = #'altdefault
    \cadenzaOn
    a'\breve g8 a8 bes4 bes4 \bar "|" \break
    a\breve bes8 a8 g4 f4 \bar "||"
    \cadenzaOff
}
\addlyrics {
    \override LyricText.font-size = #-1
    Szczęśliwy_i  \markup { \bold nie } zna kaź -- ni,
    kto_w_Pańskiej_ży -- \markup { \bold je } bo -- jaż -- ni.
}
\end{lilypond}
% 1. Szczęśliwy nie zna kaźni, \redcol{*} Kto w Pańskiej żyje bojaźni.\\
2. Najmilsza je\textbf{mu} jest droga, \redcol{*} iść według przy\textbf{ka}zań Boga.\\
3. Krew jego zac\textbf{na} na ziemi, \redcol{*} porówna zaw\textbf{żdy} z możnymi.\\
4. Ród się cnotli\textbf{wych} rozpieni, \redcol{*} i będą bło\textbf{go}sławieni.\\
5. Dom jego bę\textbf{dzie} obfity, \redcol{*} w zbiory i trwa\textbf{łe} zaszczyty.\\
6. A cześć poczci\textbf{we}go człeka, \redcol{*} pójdzie od wie\textbf{ka} do wieka.\\
7. Niechaj noc zać\textbf{mi} mgły swymi, \redcol{*} światło jest nad \textbf{cno}tliwymi.\\
8. Zawsze im Pan \textbf{jest} życzliwy, \redcol{*} litosny i \textbf{spra}wiedliwy.\\
9. Człek miły peł\textbf{en} wesela, \redcol{*} gdy drugiemu z swe\textbf{go} udziela.\\
10. W słowie się swo\textbf{im} tak rządzi, \redcol{*} że, co wyrze\textbf{cze}, nie zbłądzi.\\
11. Ludzka pamięć \textbf{jego} sprawy \redcol{*} uwieczni: „Był \textbf{to} mąż prawy”.\\
12. I zły go ję\textbf{zyk} nie trwoży, \redcol{*} bezpieczny w nad\textbf{zie}i Bożej.\\
13. Stateczność je\textbf{go} cnej duszy, \redcol{*} żadną się rze\textbf{czą} nie wzruszy.\\
14. Wydoła złej \textbf{chwi}li snadnie, \redcol{*} aż nieprzyja\textbf{ciel} upadnie.\\
15. Rozsypał swo\textbf{je} szczodroty, \redcol{*} na wdowy, bied\textbf{ne} sieroty.\\
16. Przeto uczyn\textbf{no}ścią słynie, \redcol{*} i chwała je\textbf{go} nie zginie.\\
17. Zły, na to pa\textbf{trząc}, boleje, \redcol{*} zgrzyta, z zazdro\textbf{ści} sinieje.\\
18. Taki, co w gło\textbf{wie} uradzi, \redcol{*} do skutku nie \textbf{do}prowadzi.\\
19. Chwała Ojcu \textbf{i} Synowi, \redcol{*} oraz Święte\textbf{mu} Duchowi.\\
20. Jak na począ\textbf{tku}, tak ninie, \redcol{*} i na wieki \textbf{nie}chaj słynie.\\

\proroctwo{Antyfona}

\gregorioscore{Melody/an--ecce_vir_oriens--solesmes_1961.1.gabc}

\medskip
Oto mąż, a Wschód imię Jego, zasiądzie i zapanuje i pokój narodom opowie.

\proroctwo{Antyfona}

\gregorioscore{Melody/an--dominus_judex_noster--solesmes.gabc}

\medskip
Pan Sędzią naszym, Pan Prawodawcą naszym, Pan Królem naszym. Ten zbawi nas.

\proroctwo{Psalm 112}

\begin{lilypond}
\relative c' {
    \key g \major
    \clef treble
    \override Staff.TimeSignature.stencil = ##f
    \override Staff.NoteHead.style = #'altdefault
    \cadenzaOn
    g'8 d\breve e8 fis8 g8([fis8]) e8 d8([c8]) b4 \bar "|" \break
    a'\breve b8 c8 b8 a4 g4 \bar "||"
    \cadenzaOff
}
\addlyrics {
    \override LyricText.font-size = #-1
    Chwal -- cie,_o_dziatki,  \markup { \bold Naj } -- wyż -- sze -- go Pa -- na.
    Niech_Mu_jednemu_cześć \markup { \bold bę} -- dzie śpie -- wa -- na.
}
\end{lilypond}

% 1. Chwalcie o dziatki, Najwyższego Pana, \redcol{*} niech Mu jednemu cześć będzie śpiewana.\\
2. Niech Imię Pańskie \textbf{prze}błogosławione \redcol{*} po wszystkie wieki bę\textbf{dzie} uwielbione.\\
3. Gdzie wschodzi słońce \textbf{i} kędy zapada, \redcol{*} niechaj świat Boską chwa\textbf{łę} opowiada.\\
4. Pan ma narody \textbf{wszyst}kie pod nogami. \redcol{*} Jego się chwała wzno\textbf{si} nad gwiazdami.\\
5. Kto się wżdy z Panem \textbf{tym} porówna, który \redcol{*} wysoko siedząc, z swej \textbf{nie}bieskiej góry.\\
6. I co na niebie, \textbf{i} co jest na ziemi, \redcol{*} oczyma widzi nie\textbf{u}chronionymi.\\
7. On ubogiego \textbf{z nę}dzy wyprowadzi \redcol{*} i z książętami na \textbf{ła}wie posadzi.\\
8. On niesie radość \textbf{dla} niepłodnej matki, \redcol{*} miłe w jej domu roz\textbf{mna}żając dziatki.\\
9. Chwała bądź Bogu \textbf{w Trój}cy jedynemu, \redcol{*} Ojcu, Synowi, Du\textbf{cho}wi Świętemu.\\
10. Jak od początku \textbf{by}ła, tak i ninie, \redcol{*} I na wiek wieków nie\textbf{chaj} zawsze słynie.\\

\proroctwo{Antyfona}

\gregorioscore{Melody/an--dominus_judex_noster--solesmes.gabc}

\medskip
Pan Sędzią naszym, Pan Prawodawcą naszym, Pan Królem naszym. Ten zbawi nas.

\proroctwo{Antyfona}

\gregorioscore{Melody/an--ecce_dedi_te--solesmes_1961.1.gabc}

\medskip
Oto uczyniłem Cię światłością narodów, abyś był zbawieniem mojem aż na krańce ziemi.

\proroctwo{Psalm 116}

\begin{lilypond}
\relative c' {
    \key d \minor
    \clef treble
    \override Staff.TimeSignature.stencil = ##f
    \override Staff.NoteHead.style = #'altdefault
    \cadenzaOn
    c'8 bes8 a8 a8 a8 a8 a8 bes8 a8 g4 f4 \bar "|" \break
    g\breve a8([g8]) f8 e8 d4 \bar "||"
    % bes'8 aes8 g\breve aes8 g8 f4 ees4 \bar "|" \break
    % f\breve g8([f8]) ees8 d4 c4 \bar "||"
    \cadenzaOff
}
\addlyrics {
    \override LyricText.font-size = #-1
    Bo -- ga na -- sze -- go chwal -- cie \markup { \bold wszyst } -- kie zie -- mie,
    Daj_mu_cześć_winną_całe  \markup { \bold ludz } -- kie ple -- mię.
}
\end{lilypond}

% 1. Boga naszego chwalcie \textbf{wszyst} kie ziemie, \redcol{*} Daj mu cześć winną całe \textbf{lu}dzkie plemię.\\
2. Bo litość jego nad na\textbf{mi} stwierdzona, \redcol{*} A prawda Pańska wiecznie uiszczona.\\
3. Chwała bądź Bogu w Trójcy \textbf{je}dynemu, \redcol{*} Ojcu, Synowi, Ducho\textbf{wi} Świętemu.\\
4. Jak od początku była \textbf{tak} i ninie, \redcol{*} i na wiek wieków niechaj \textbf{zaw}sze słynie.\\

\proroctwo{Antyfona}

\gregorioscore{Melody/an--ecce_dedi_te--solesmes_1961.1.gabc}

\medskip
Oto uczyniłem Cię światłością narodów, abyś był zbawieniem mojem aż na krańce ziemi.

\proroctwosigl{Kapitulum}{Kol 1:12-13}

\oremuss{
    Fratres: Grátias ágimus Deo Patri, 
    qui dignos nos fecit in partem sortis sanctórum in lúmine, 
    qui erípuit nos de potestáte tenebrárum, 
    et tránstulit in regnum Fílii dilectiónis suæ.
}{
    Bracia: dzięki czynimy Bogu Ojcu, który nas godnymi uczynił dziedzictwa Świętych w światłości; 
    który nas wyrwał z mocy ciemności i przeniósł w królestwo Syna umiłowanego Swego.
}
\oremuss{
    \rr Deo grátias.
}{
    \rr Bogu niech będą dzięki.
}

\proroctwo{Hymn}

% \grechangestaffsize{15}
\gregorioscore{Melody/hy--te_saeculorum_principem--solesmes.gabc}
% \grechangestaffsize{17}

\bigskip
\oremuss{
    \textbf{1.} Ciebie, Książeciem wszystkich wieków,\\
    Ciebie, narodów, Królem Chryste,\\
    Ciebie też dusz i serc jedynym,\\
    Wszechwładnym wyznawamy Panem.\\

    \textbf{2.} Zbrodnicza tłuszcza woła:\\
    Nie chcem, ażeby nami Chrystus władał!\\
    My zaś obwołujemy, Panie,\\
    Ciebie, najwyższym Królem wszystkich.\\

    \textbf{3.} Książę pokoju Chryste, poddaj\\
    Sobie umysły buntownicze,\\
    A owce, co od Twej miłości\\
    Odbiegły, zgromadź do owczarni.\\

    \textbf{4.} Wszak potoś wisiał sam na krzyżu,\\
    Potoś otworzył swe ramiona,\\
    I włócznią ranę Ci zadano,\\
    W której k'nam serce Twoje płonie
}{
    \textbf{5.} Dlatego kryjesz się w ołtarzu\\
    Pod postaciami wina, chleba,\\
    I z przebitego boku swego\\
    Zlewasz na dzieci wybawienie.\\

    \textbf{6.} Niechaj narodów wodze Tobie\\
    Publiczne hołdy i cześć niosą,\\
    Nauczyciele, sędzie, prawo\\
    I sztuka niech Cię wyrażają.\\

    \textbf{7.} Niechaj oznaki władzy królów\\
    Tem błyszczą, że oddane Tobie,\\
    Rządź słodkiem berłem Twem Ojczyzną,\\
    Obywateli rządź domami.\\

    \textbf{8.} Jezu, niech Tobie będzie chwała,\\
    Który berłami świata władasz,\\
    Z Ojcem i Duchem Świętym, Boże,\\
    Przez wszystkie wieki wieków.\\
    Amen.
}

\bigskip

\gresetinitiallines{0}
\gregorioscore{Melody/hy--te_saeculorum_principem--solesmes_1961.2.gabc}
\gregorioscore{Melody/hy--te_saeculorum_principem--solesmes_1961.3.gabc}
\gresetinitiallines{1}

% \rr Et pacis non erit fi\textbf{nis}.

\oremuss{}{
\vv Panowanie Jego rozszerzy się.\\
\rr A pokojowi nie będzie końca.
}

\newpage
\proroctwo{Magnificat}

\gregorioscore{Melody/an--habet_in_vestimento--solesmes_1961.1.gabc}

\medskip
Na szacie i na biodrze swojem ma napisane Król królów i Pan panujących. Jemu chwała i władza na wieki wieków.

\bigskip

\begin{lilypond}
\relative c' {
    \key f \major
    \clef treble
    \override Staff.TimeSignature.stencil = ##f
    \override Staff.NoteHead.style = #'altdefault
    \cadenzaOn
    f\breve g8 g8 a8([g8]) f4 \bar "|" \break
    f\breve e8 f8 d4 c4 \bar "|" \break
    a'\breve g8 a8 bes4 bes4 \bar "|" \break
    a\breve g8 f8 g4 f4 \bar "||"
    \cadenzaOff
}
\addlyrics {
    \override LyricText.font-size = #-1
    Uwielbiaj,_duszo_moja,_sławę \markup { \bold Pa } -- na me -- go,
    chwal_Boga_Stworzyciela_tak_bar -- \markup { \bold dzo } dob -- re -- go.
    Bóg_mój,_zbawienie_moje,_jedy -- \markup { \bold na } o -- tu -- cha,
    Bóg_mi_rozkoszą_serca_i_we -- \markup { \bold se } -- lem du -- cha.
}
\end{lilypond}

% 1. Uwielbiaj, \ding{64} duszo moja, sławę \textbf{Pa}na mego, \redcol{*} Chwal Boga Stworzyciela tak bar\textbf{dzo} dobrego.\\
% 2. Bóg mój, zbawienie moje, jedy\textbf{na} otucha, \redcol{*} Bóg mi rozkoszą serca i we\textbf{se}lem ducha.\\
3. Bo mile przyjąć raczył swej słu\textbf{gi} pokorę, \redcol{*} łaskawym okiem wejrzał na Da\textbf{wi}da córę.\\
4. Przeto wszystkie narody, co zie\textbf{mię} osiędą, \redcol{*} Odtąd błogosławioną mnie na\textbf{zy}wać będą.\\
5. Bo wielkimi darami uczczo\textbf{nam} od Tego, \redcol{*} którego moc przedziwna, święte \textbf{I}mię Jego.\\
6. Którzy się Pana boją, szczęśli\textbf{wi} na wieki, \redcol{*} bo z nimi miłosierdzie z rodu \textbf{w ród} daleki.\\
7. Na cały świat pokazał moc swych \textbf{ra}mion świętych, \redcol{*} rozproszył dumne myśli głów py\textbf{chą} nadętych.\\
8. Wyniosłych złożył z tronu, znikczem\textbf{nił} wielmożne, \redcol{*} wywyższył, uwielmożnił w poko\textbf{rę} zamożne.\\
9. Głodnych nasycił, hojnie i w dob\textbf{ra} spanoszył, \redcol{*} bogaczów z torbą puścił i nędz\textbf{nie} rozproszył.\\
10. Przyjął do łaski sługę Izra\textbf{e}la cnego, \redcol{*} wspomniał nań, użyczył mu miło\textbf{sier}dzia swego.\\
11. Wypełnił, co był przyrzekł niegdyś \textbf{oj}com naszym: \redcol{*} Abrahamowi z potomstwem, jego \textbf{wiecz}nym czasem.\\
{\footnotesize\color{red}{\textbf{Improwizacja organowa do czasu, aż turyfer powróci od okadzenia ludu.}}}\\
12. Wszyscy śpiewajmy Bogu w Trójcy \textbf{Je}dynemu, \redcol{*} Chwała Ojcu, Synowi, Ducho\textbf{wi} Świętemu.\\
13. Jak była na początku, tak zaw\textbf{sze} niech będzie, \redcol{*} teraz i na wiek wieków niechaj \textbf{sły}nie wszędzie.\\

\gregorioscore{Melody/an--habet_in_vestimento--solesmes_1961.1.gabc}

\medskip
Na szacie i na biodrze swojem ma napisane Król królów i Pan panujących. Jemu chwała i władza na wieki wieków.

\proroctwomedskip{Oracja}

\oremuss{
    \vv Dominus vobiscum.
}{
    \vv Pan z wami.
}
\oremuss{
    \rr Et cum spiritu tuo.
}{
    \rr I z duchem twoim.
}
\oremuss{
    Orémus. Omnípotens sempitérne Deus, qui in dilécto Fílio tuo, universórum Rege, ómnia instauráre voluísti:
    concéde propítius; ut cunctæ famíliæ géntium, peccáti vúlnere disgregátæ, eius suavíssimo subdántur império:
    Qui tecum vivit et regnat in unitáte Spíritus Sancti, Deus, per ómnia sǽcula sæculórum.
}{
    Módlmy się. Wszechmogący wieczny Boże, który przez ukochanego Syna swego, Króla wszechświata, 
    wszystko odnowić postanowiłeś spraw miłościwie, by wszystkie rodziny narodów, 
    wskutek ran grzechowych rozdzielone, poddały się Jego najsłodszemu Panowaniu.
    Który z Tobą żyje i króluje w jedności Ducha Świętego Bóg przez wszystkie wieki wieków.
}
\oremuss{
    \rr Amen.
}{
    \rr Amen.
}


\proroctwo{Zakończenie}

\oremuss{
    \vv Dominus vobiscum.
}{
    \vv Pan z wami.
}
\oremuss{
    \rr Et cum spiritu tuo.
}{
    \rr I z duchem twoim.
}

% \vv Benedicamus Domino.

% \rr Deo gratias.

\gresetinitiallines{0}
% \grecommentary[0.1cm]{Orbis Factor (XI)}
\gregorioscore{Melody/or--benedicamus_domino_i_classis_in_ii_vesperis_mode_5--solesmes_1961.1.gabc}
\gresetinitiallines{1}

\medskip
\oremuss{}{
\vv Błogosławmy Panu.\\
\rr Bogu niech będą dzięki.
}
\oremuss{
    \vv Fidelium animae per misericordiam Dei requiescant in pace.
}{
    \rr Dusze wiernych zmarłych przez miłosierdzie Boże niech odpoczywają w pokoju.
}
\oremuss{
    \rr Amen.
}{
    \rr Amen.
}


\proroctwo{Antyfona}

\gregorioscore{Melody/an--salve_regina_simplex-solesmes.1.gabc}

\vfill

\centerline{\pgfornament[width=3cm]{162}}

\vfill

\newpage

\thispagestyle{empty}

\vfill

\begin{tikzpicture}[color=black,
every node/.style={inner sep=0pt}]
\node[minimum width=12.8cm, minimum height=18.0cm] (vecbox) at (current page.center) {};
\node[align=center, text width=10cm, inner sep=10pt] at (vecbox.center) {%
    {\Large \textbf{A. M. D. G.}}\\[1em]
};
\end{tikzpicture}

\vfill

\end{document}
